% (C) 2026 Ralf Stephan, in collaboration with Claude Code.
% Released under CC0 1.0 Universal (public-domain dedication).
% See https://creativecommons.org/publicdomain/zero/1.0/

\documentclass[11pt]{article}
\usepackage[T1]{fontenc}
\usepackage[utf8]{inputenc}
\usepackage{lmodern}
\usepackage[margin=2.5cm]{geometry}
\usepackage{amsmath,amssymb}
\usepackage{braket}
\usepackage[hidelinks]{hyperref}

% Lean names and paths: verbatim, breakable at . _ /
\DeclareUrlCommand\lean{\urlstyle{tt}}
\def\UrlBreaks{\do\.\do\_\do\/}
% Extra stretch for paragraphs with long, unbreakable Lean names
\emergencystretch=6em

\title{Open problems:\\ Entropy is not continuous in standard probability spaces}
\author{}
\date{}

\begin{document}
\maketitle

\noindent
Open problems raised by the technical brief in this directory, and by its Lean formalization
(see \lean{formalization-log.tex}).

\section*{From the brief}
\begin{enumerate}
\item \textbf{The appropriate topology on $P_\mu$.} The brief raises ``the question as to what is
  the appropriate one for probability spaces'' and shows only that a topology making mixing and
  $S$ continuous must refine total variation. Whether such a topology exists with good properties
  (e.g.\ Hausdorff), and which one is the coarsest, is left open. (Brief, Sec.~1; Theorem~46;
  Sec.~5.)
\item \textbf{When is entropy $L^1$-continuous?} Theorem~21 needs disjoint sets of arbitrarily
  large finite measure; the conclusion extends the failure to ``typical infinite-dimensional
  spaces''. Measure spaces outside the hypothesis, such as Lebesgue measure on $[0,1]$, are not
  treated, and no characterization of the spaces on which $S$ is $L^1$-continuous is given.
  (Brief, Theorem~21; Sec.~5.)
\end{enumerate}

\end{document}
